% Imported from https://github.com/Luca-Yucheng-Jin/QFT-soln
% Source commit: 0bb3e7066154a8fdb256399a8ecceae720c5e192
\section{David Tong The Standard Model, Sheet 2: Symmetry Breaking and QCD}
\markboth{\thesection\quad TONG THE STANDARD MODEL, SHEET 2}{}

\noindent\textit{These prompts paraphrase David Tong's \emph{The Standard Model},
Example Sheet 2 (February 2026). Problem numbers, starred status, equations,
and guidance follow the
\href{https://davidtong.org/pdfs/teaching/standard-model/rw2.pdf}{official sheet}.}

\subsection{Sheet 2, Problem 1: Broken Discrete Symmetry}

Construct a polynomial potential for a complex scalar $\phi$ whose
$\mathbb Z_N$ symmetry breaks spontaneously. Identify every ground
state.
\begin{solution}
    The potential
    \begin{equation}
        V(\phi) = \lambda|\phi|^{2N} - \frac{1}{2}\mu^2 (\phi^N + \phi^{*N})
    \end{equation}
    is invariant under $\mathbb{Z}_N: \phi \to e^{2\pi ik/N}\phi$. Therefore, the vacuum is given by
    \begin{equation}
        v_k = \left(\frac{\mu^2}{2\lambda}\right)^{1/N}e^{2\pi i k /N}.
    \end{equation}
\end{solution}

\subsection{Sheet 2, Problem 2: Matrix Order Parameters and Goldstone Modes}

Let $M(x)$ be a complex $N\times N$ matrix with $\lambda>0$ and
\begin{equation}
    S=\int d^4x\,\operatorname{Tr}\left[
    \partial^\mu M^\dagger\partial_\mu M
    -kM^\dagger M-\frac\lambda2M^\dagger M M^\dagger M\right].
\end{equation}
\begin{enumerate}[label=\roman*)]
    \item Show invariance under $M\mapsto AMB^\dagger$ for
    $A,B\in U(N)$ and identify the $U(1)$ subgroup of
    $U(N)\times U(N)$ that acts trivially.
    \item When $k<0$, show that a vacuum satisfies
    $M_0^\dagger M_0=v^2\mathbf1$. Find the unbroken group,
    describe the vacuum manifold as a coset, and count Goldstone bosons.
    \item Add
    $\int d^4x\,h(\det M+\det M^\dagger)$.
    Determine the new symmetry group and the Goldstone count if the
    vacuum still has $M_0^\dagger M_0=v^2\mathbf1$ with $v\ne0$.
\end{enumerate}
\begin{solution}
    \begin{enumerate}[label=\roman*)]
    \item  The invariance is trivial to see. For $A = B \in U(1)$, the transformation is trivial.
    \item The differentiation of the potential with respect to $M^\dagger$, we find
    \begin{equation}
        \frac{\partial}{\partial M^\dagger}\left(\frac{\lambda}{2}M_0^\dagger M_0 M_0^\dagger M_0 + k M_0^\dagger M_0\right) = M_0 (\lambda M_0^\dagger M_0 + k ) = 0 \implies M_0^\dagger M_0 = -\frac{k}{\lambda}\mathbf{1}.
    \end{equation}
    Therefore, the vacuum must take the form $M_0 = v U$ with $U \in U(N)$. We can choose $U = \mathbf{1}$.
    Under $U(N) \times U(N)$, the vacuum transforms as
    \begin{equation}
        v\mathbf{1} \to  A v\mathbf{1} B^\dagger.
    \end{equation}
    Therefore, the vacuum is invariant for $A = B \in U(N)$ . Therefore, the stability group is given by $H = U(N)$, while the coset space is $G = (U(N)\times U(N))/U(N) $. Therefore, the number of Goldstone bosons is $\dim G = N^2$.
    \item With this additional term, the symmetry reduces to $A, B \in SU(N)$, namely $SU(N) \times SU(N)$. If the vacuum still takes the same form, then the stability group becomes $H = SU(N)$. Hence, the number of Goldstone bosons becomes $\dim SU(N) = N^2 -1$.
    \end{enumerate}
\end{solution}

\subsection{Sheet 2, Problem 3*: Fundamental Higgs Fields in $SU(2)$ and $SU(N)$}

For a fundamental $SU(2)$ scalar $\phi_a$ ($a=1,2$), use
\begin{equation}
    S=\int d^4x\left[-\frac12\operatorname{Tr}(F_{\mu\nu}F^{\mu\nu})
    +D_\mu\phi^\dagger D^\mu\phi
    -\frac\lambda2(\phi^\dagger\phi-v^2)^2\right],
\end{equation}
where $F_{\mu\nu}=\partial_\mu A_\nu-\partial_\nu A_\mu
-ig[A_\mu,A_\nu]$, $D_\mu\phi=\partial_\mu\phi-igA_\mu\phi$,
and $T^a=\sigma^a/2$ in the fundamental representation.
Find the particle masses. Generalise to a fundamental scalar in
$SU(N)$: after condensation, what symmetry survives, how many gauge
bosons become massive, and how many real scalar components are eaten?
\begin{solution}
    For $SU(2)$, we can choose the vacuum to be
    \begin{equation}
        \phi_0 = \begin{pmatrix}
            0 \\ v
        \end{pmatrix}.
    \end{equation}
    Therefore, all of the symmetries are broken, and we can reparameterise the scalar field as
    \begin{equation}
        \phi(x) = \begin{pmatrix}
            0 \\ v + \sigma(x)
        \end{pmatrix}e^{i\pi^a(x)T^a/ F_\pi},
    \end{equation}
    where $F_\pi $ is chosen to by $v$ so that the kinetic term of $\pi$ is canonically normalised. The gauge boson acquires a mass through the term
    \begin{equation}
        g^2A_\mu A^\mu |\phi|^2 = v^2g^2A_\mu A^\mu + \cdots,
    \end{equation}
    where $m_A = \sqrt{2}vg$. The $\pi(x)$ field is eaten by the gauge bosons. The mass of $\sigma(x)$ is given by
    \begin{equation}
        \lambda v^2 |\phi|^2 = \lambda v^2 \sigma^2 + \cdots,
    \end{equation}
    which gives $m_\sigma = \sqrt{2\lambda}v$.

    For $SU(N)$, a choice of vacuum could be
    \begin{equation}
        \phi_0 = \begin{pmatrix}
            0 \\ \vdots \\ v
        \end{pmatrix}.
    \end{equation}
    Therefore, the stability group is $H = SU(N-1)$, and hence the coset space is $G = SU(N) / SU(N-1)$. Therefore, we can reparameterise the scalar field as
        \begin{equation}
        \phi(x) = mm\begin{pmatrix}
            0 \\\vdots \\ v + \sigma(x)
        \end{pmatrix}e^{i\pi^a(x)T^a/ F_\pi},
    \end{equation}
    where $T^a$ are the generators of $G$ and $a = 1, \cdots , \dim G$. For consistency, $\phi(x)$ has $2N$ degrees of freedom, while $\pi(x)$ has $\dim G = 2N-1$ degrees of freedom and $\sigma(x)$ has $1$ degree of freedom. All of the $\pi$ are eaten by the gauge fields. Therefore, there will be $2N-1$ massive bosons.
\end{solution}

\subsection{Sheet 2, Problem 4: Adjoint Higgs Breaking Patterns}

An $SU(N)$ adjoint scalar is a traceless matrix
$\phi=\phi^AT^A$ with
$D_\mu\phi=\partial_\mu\phi-ig[A_\mu,\phi]$.
Explain why a potential minimum can be represented by
\begin{equation}
    \phi_0=\operatorname{diag}(v_1,\ldots,v_N),
    \qquad \sum_{a=1}^Nv_a=0,\quad v_a\leq v_{a+1}.
\end{equation}
Classify the residual gauge groups according to coincidences among
the $v_a$.
\begin{solution}
    Given a vacuum $\phi_0$, we can always diagonalise it by $U\phi_0 U^\dagger$, where $U \in SU(N)$. Therefore, we can choose the vacuum to be
    \begin{equation}
        \phi_0  = \operatorname{diag}(v_1,\ldots,v_N),\quad v_a\leq v_{a+1}.
    \end{equation}
    Since $T^A$ are traceless, we must have $\Tr \phi_0 = 0$, which implies $\sum_a v_a = 0$. For a gauge transformation generated by $T^A$ to be unbroken, we must have
    \begin{equation}
        [T^A, \phi_0] = 0.
    \end{equation}
    In general, we have
    \begin{equation}
        [T^A,\phi_0]_{ik} = \sum_j T^A_{ij}\delta_{jk}v_k - \sum_j \delta_{ij}v_j T^A_{jk} = (v_k -v_i) T^A_{ik}.
    \end{equation}
    Therefore, there are no restrictions on $T_{ii}$, which gives $N-1$ unbroken generators. Meanwhile, we must either have $T_{ik} = 0$ for all $i\neq k$, or $v_i - v_k = 0$ for some $i,k$. Therefore, the number of unbroken generator is equal to the number of identical eigenvalue pairs. Suppose each distinct eigenvalues have multiplicity $n_i$, they give
    \begin{equation}
        2\sum_{j}\begin{pmatrix}
            n_j \\ 2
        \end{pmatrix} = \sum_j(n_j^2 -n_j) = \sum_jn_j^2 - N.
    \end{equation}
    Therefore, the total number of unbroken generators is
    \begin{equation}
        \sum_jn_j^2 - 1.
    \end{equation}
\end{solution}

\subsection{Sheet 2, Problem 5*: One-Loop QCD Scale}

For $SU(N_c)$ with $N_f$ massless fundamental Dirac fermions, use
\begin{equation}
    \frac1{g^2(\mu)}=\frac1{g_0^2}
    -\frac{11N_c-2N_f}{3(4\pi)^2}
    \log\frac{\Lambda_{\rm UV}^2}{\mu^2}.
\end{equation}
Taking $\alpha_s=g_s^2/(4\pi)\simeq0.12$ at
$\mu=M_Z\simeq90\,\mathrm{GeV}$, estimate $\Lambda_{\rm QCD}$
if every quark lighter than $M_Z$ is treated as massless.
How can the estimate be made more realistic?
\begin{solution}
    Taking the physical scale, we have
    \begin{equation}
        \frac{1}{g^2(M_Z)} =  \frac{1}{g_0^2} -\frac{11N_c-2N_f}{3(4\pi)^2}
    \log\frac{\Lambda_{\rm UV}^2}{M_Z^2},
    \end{equation}
    while the Landau pole is at
    \begin{equation}
         0 = \frac{1}{g^2(\Lambda_\mathrm{QCD})} =  \frac{1}{g_0^2} -\frac{11N_c-2N_f}{3(4\pi)^2}
    \log\frac{\Lambda_{\rm UV}^2}{\Lambda_\mathrm{QCD}^2}.
    \end{equation}
    Combining the two, we find
    \begin{equation}
        \frac{11N_c-2N_f}{3(4\pi)^2}
    \log\frac{M_Z^2}{\Lambda_\mathrm{QCD}^2} = \frac{1}{g^2(M_Z)} \implies \Lambda_{\mathrm{QCD}}=  M_Z\exp\left(-\frac{6\pi}{(11N_c - 2N_f)\alpha_s}\right).
    \label{eq:LandauPole}
    \end{equation}
    Five quark flavours are lighter than \(M_Z\): \(u,d,s,c,b\). The top quark is heavier, so take \(N_c=3\) and \(N_f=5\). This gives $\Lambda_\mathrm{QCD} \approx 0.097\,\mathrm{GeV}$.

    For a more realistic estimation, we can run the coupling towards the IR. When we go across the boundary of the bottom quark mass, we take $N_f$ to be $4$. Then, we match the value at $\alpha(m_b)$ in both regime. We repeat this procedure until ran into the landau pole.
\end{solution}

\subsection{Sheet 2, Problem 6: Casimirs and Dynkin Indices}

For a representation $R$, define
\begin{equation}
    T^AT^A=C(R)\mathbf1,
    \qquad
    \operatorname{Tr}(T^AT^B)=\frac12I(R)\delta^{AB}.
\end{equation}
Prove $2C(R)\dim(R)=I(R)\dim(G)$.
Find $C(R)$ for the fundamental and anti-fundamental of $SU(N)$,
and calculate the adjoint Casimir of $SU(2)$.
\begin{solution}
    \begin{equation}
        \Tr(T^A T^A) = C(R) \dim \mathbf{1} = C(R)\dim R,
    \end{equation}
    while
    \begin{equation}
        \Tr(T^AT^A) = \frac{1}{2}I(R) \delta^{AA} = \frac{1}{2}I(R) \dim G.
    \end{equation}
    Therefore,
    \begin{equation}
        2C(R)\dim(R)=I(R)\dim(G).
    \end{equation}
    For fundamental representation of $SU(N)$, we have $\dim G = \dim \mathbf{N}$ and $I(\mathbf{N}) =1$, and hence $C(\mathbf{N})  = \frac{1}{2}$. Since $\bar T^A = - (T^A)^*$, we must have $I(\mathbf{N}) = I(\bar{\mathbf{N}})$, and hence $C(\mathbf{N}) = C(\bar{\mathbf{N}}) = 1/2$.

    For adjoint representation of $SU(2)$, the generators are$-i\epsilon^i{}_{jk}$, and hence
    \begin{equation}
        i\epsilon^i{}_{jk}i\epsilon^i{}_{kl} =2\delta_{jl} \implies C(\mathrm{adj}) = 2.
    \end{equation}
\end{solution}

\subsection{Sheet 2, Problem 7: Expanding the Chiral Lagrangian}

For $U=e^{2i\pi/f_\pi}$ with $\pi\in su(N_f)$, expand
\begin{equation}
    \mathcal L_{\rm pion}
    =\frac{f_\pi^2}{4}
    \operatorname{tr}(\partial^\mu U^\dagger\partial_\mu U)
\end{equation}
through fourth order in $\pi$. Show that the quadratic term is
$\operatorname{tr}(\partial_\mu\pi)^2$ and that the quartic term
can be written
\begin{equation}
    -\frac{2}{3f_\pi^2}\operatorname{tr}\left[
    \pi^2(\partial_\mu\pi)^2
    -(\pi\partial_\mu\pi)^2\right].
\end{equation}
For $N_f=2$ and $T^a=\sigma^a/2$, reduce the interaction to
\begin{equation}
    \mathcal L_{\rm int}=-\frac1{6f_\pi^2}
    \left[\pi^a\pi^a\partial_\mu\pi^b\partial^\mu\pi^b
    -\pi^a\partial_\mu\pi^a\,
    \pi^b\partial^\mu\pi^b\right].
\end{equation}
\begin{solution}
    Expanding $U$, we find
    \begin{equation}
        U = 1+ \frac{2i}{f_\pi}\pi - \frac{2}{f_\pi^2}\pi^2 - \frac{4i}{3f_\pi^3}\pi^3 + \cdots,
    \end{equation}
    and hence
    \begin{equation}
\partial_\mu U
=\frac{2i}{f_\pi}\partial_\mu\pi
-\frac{2}{f_\pi^2}\bigl((\partial_\mu\pi)\pi+\pi(\partial_\mu\pi)\bigr)
-\frac{4i}{3f_\pi^3}\bigl((\partial_\mu\pi)\pi^2+\pi(\partial_\mu\pi)\pi+\pi^2(\partial_\mu\pi)\bigr)
+\cdots.
    \end{equation}
    Expanding the Lagrangian, we have
    \begin{equation}
        \begin{aligned}
            \mathcal L_{\rm pion}
            &=\operatorname{tr}(\partial_\mu \pi\partial^\mu \pi)-\frac{2}{3f_\pi^2}\operatorname{tr}\left[
    \pi^2(\partial_\mu\pi)^2
    -(\pi\partial_\mu\pi)^2\right] +\cdots.
        \end{aligned}
    \end{equation}

    For $N_f = 2$, we need to evaluate
    \begin{equation}
        \operatorname{tr}(\sigma^a \sigma^b \sigma^c \sigma^d ) = \operatorname{tr}(\sigma^a \sigma^b \{\sigma^c, \sigma^d \}) - \operatorname{tr}(\sigma^a \{\sigma^b, \sigma^d\} \sigma^c )  + \operatorname{tr}(\{\sigma^a ,\sigma^d\} \sigma^b \sigma^c ) + \operatorname{tr}( \sigma^d \sigma^a \sigma^b \sigma^c).
    \end{equation}
    Using $\{\sigma^a,\sigma^b\} = 2\delta^{ab}$ and $\Tr{\sigma^a \sigma^b} = 2\delta^{ab}$, we find
    \begin{equation}
        \operatorname{tr}(\sigma^a \sigma^b \sigma^c \sigma^d ) = 2(\delta^{ab}\delta^{cd}-\delta^{bd}\delta^{ac} + \delta^{ad}\delta^{bc}).
    \end{equation}
    Substituting into the Lagrangian, the interaction reduces to
    \begin{equation}
    \mathcal L_{\rm int}=-\frac1{6f_\pi^2}
    \left[\pi^a\pi^a\partial_\mu\pi^b\partial^\mu\pi^b
    -\pi^a\partial_\mu\pi^a\,
    \pi^b\partial^\mu\pi^b\right].
\end{equation}
\end{solution}

\subsection{Sheet 2, Problem 8: Meson Mass Relations}

For $N_f=3$, use the pion, kaon, and eta matrix
\begin{equation}
    \pi=\frac1{\sqrt2}
    \begin{pmatrix}
    \pi^0/\sqrt2+\eta/\sqrt6&\pi^+&K^+\\
    \pi^-&-\pi^0/\sqrt2+\eta/\sqrt6&K^0\\
    K^-&\bar K^0&-2\eta/\sqrt6
    \end{pmatrix}.
\end{equation}
With $M=\operatorname{diag}(m_u,m_d,m_s)$ and a mass term
proportional to $f_\pi\operatorname{tr}[(M+M^\dagger)\pi^2]$,
show
\begin{equation}
    \frac{m_{K^+}^2-m_{K^0}^2}{m_\pi^2}
    =\frac{m_u-m_d}{m_u+m_d}.
\end{equation}
Then take $m_u\simeq m_d$ to derive the Gell-Mann--Okubo relation
$4m_K^2\simeq3m_\eta^2+m_\pi^2$ and compare it with measured
meson masses.
\begin{solution}
\[
(\pi^2)_{11}=\frac12\left[\left(\frac{\pi^0}{\sqrt2}+\frac{\eta}{\sqrt6}\right)^2+\pi^+\pi^-+K^+K^-\right],
\]\[
(\pi^2)_{22}=\frac12\left[\left(-\frac{\pi^0}{\sqrt2}+\frac{\eta}{\sqrt6}\right)^2+\pi^-\pi^++K^0\bar K^0\right],
\]\[
(\pi^2)_{33}=\frac12\left[K^-K^++\bar K^0K^0+\frac{2\eta^2}{3}\right].
\]
Hence, the mass term is proportional to
\begin{equation}
     m_u (\pi^2)_{11} +  m_d (\pi^2)_{22} + m_s (\pi^2)_{33}.
\end{equation}
Then, we can read off the masses of the Mesons, which are
\begin{equation}
    m_{\pi}^2 =  m_u + m_d, \quad m_{K^0}^2 = m_d + m_s, \quad m_{K^+}^2 = m_u + m_s,\quad m_\eta^2 = \frac{m_u + m_d + 4m_s}{3},
\end{equation}
up to an overall constant. Therefore, we obtain the relation
\begin{equation}
    \frac{m_{K^+}^2-m_{K^0}^2}{m_\pi^2}
    =\frac{m_u-m_d}{m_u+m_d}.
\end{equation}
Take $m_u \simeq m_d$, we immediately find $m_{K}^2 \simeq m_{K^0}^2 \simeq m_{K^+}^2\simeq m_u + m_s$. Then, it follows
\begin{equation}
    4m_K^2\simeq3m_\eta^2+m_\pi^2.
\end{equation}
Using the PDG measured masses and averaging the charged and neutral squared masses within each pion and kaon multiplet gives
\[
4m_K^2 \simeq 0.983\,\mathrm{GeV}^2.
\]\[
3m_\eta^2+m_\pi^2 \simeq 0.920\,\mathrm{GeV}^2.
\]
\end{solution}

